\paragraph{Aggregate case study: relation memory and temporal context.}
Table~\ref{tab:case-relation-temporal-context} provides two controlled aggregate views of the dynamic context mechanisms in DHC. We report aggregate test metrics rather than individual edge-level rank trajectories, because the completed artifacts contain run-level metrics, train histories, and memory traces, but no exported per-edge before/after ranking samples. The relation-memory case uses the controlled \texttt{best\_config\_module\_ablation} runs for \texttt{thgl\_software\_edge\_inductive}. Removing the relation-memory state reduces AP from 0.5661 to 0.5466, a drop of 0.0195 (3.45\% relative), while MRR changes only slightly from 0.7971 to 0.7967. This pattern suggests that the relation memory is more visible in aggregate positive-edge discrimination than in the top-ranked reciprocal-rank signal alone. The effect is consistent with the role of \texttt{NodeRelationMemory}: it accumulates relation-conditioned node states over time, complementing the short-range neighbor evidence used by the batch encoder. The larger AP loss after removing the community module (0.0382) further indicates that long-range structural context and relation-specific temporal state are complementary rather than redundant. The corresponding histories and memory artifacts are stored under \texttt{results/best\_config\_module\_ablation/thgl\_software\_edge\_inductive\_\_*}.

The temporal-context sweep in Panel~(b) shows that the recent-neighbor window is also not a monotonic capacity knob. For \texttt{patent\_cn} inductive prediction, the best AP is obtained at $K=16$ (0.9870), whereas expanding the window to $K=32$ lowers AP to 0.9822. In the transductive split, the best setting shifts to $K=8$ (0.9873), and the weakest setting is $K=12$ (0.9825). These results support the interpretation that temporal context acts as a recency filter: a moderate or short history can preserve actionable recent interactions, while an overly long window may introduce stale neighbors that dilute the signal for edge prediction.
